\section{Related Work}
\label{sec:related}
Prior work relates agent activity to operating-system events for observability and provenance, resource control and enforcement, general-purpose instrumentation, and request accounting. We review it in this order and then compare it along the dimensions that network attribution requires.

AgentSight~\cite{zheng2025agentsight} correlates intercepted LLM traffic with kernel events through process lineage, a 100--500\,ms temporal window, and argument matching. Its network events are per-process summaries, and its evaluation reports overhead and case studies, not correlation accuracy. AgentProvenance~\cite{byteyellow2026agentprovenance} is the closest system: it resolves syscall-level \texttt{connect} events to execution-scope and tool-call bindings through cgroup, process, and time tiers, and it marks unmatched events as \emph{unresolved}. However, its events carry no connection identity, overlapping bindings are resolved silently (latest start wins) rather than reported as ambiguous, and its validation consists of pass/fail acceptance tests. King et al.~\cite{king2026kernelevidence} validate kernel evidence at session level and state explicitly that per-tool-call attribution remains an open problem.

AgentCgroup~\cite{zheng2026agentcgroup} places each shell tool call in its own cgroup for CPU and memory control; it does not account network activity. ActPlane~\cite{zheng2026actplane} and AgentGuard~\cite{shukla2026agentguard} enforce network policies over process-tree domains or tracked PIDs; AgentGuard keys denied connections to a single globally active session. ContractWarden~\cite{cui2026contractwarden} stores a task identifier in BPF task-local storage, inherited across fork and exec, and enforces policies at \texttt{socket\_connect}; its unit is the launched task, not the individual tool call, and it does not report ambiguity or validate attribution. IsolateGPT~\cite{wu2025isolategpt} isolates the apps of an agent system in separate execution environments, and Progent~\cite{shi2025progent} enforces privilege policies on tool calls inside the agent; neither attributes the resulting kernel network activity.

OpenTelemetry eBPF Instrumentation (OBI)~\cite{otel2026obi} produces HTTP/gRPC spans with W3C trace-context propagation across threads and event loops, and exports flow metrics labeled with workload attributes. In the inspected source, context does not follow forked children and flow metrics carry no request identity. DeepFlow~\cite{shen2023deepflow} assembles distributed traces from eBPF-captured network messages without code changes, and Tetragon~\cite{tetragon} enforces process-aware policies in the kernel; both key activity to processes, containers, or protocol-level requests, not to agent executions. The OpenTelemetry GenAI conventions define a tool-call identifier at the application level~\cite{otel2026genai}, which kernel network events do not carry.

Charging kernel activity to logical requests predates agents: resource containers decouple accounting from processes~\cite{banga1999resource}; Magpie and Whodunit reconstruct requests from kernel and inter-process events~\cite{barham2004magpie,chanda2007whodunit}; X-Trace, Dapper, Canopy, Pivot Tracing, and Retro propagate request context for tracing and per-workflow resource management~\cite{fonseca2007xtrace,sigelman2010dapper,kaldor2017canopy,mace2015pivot,mace2015retro}; and provenance systems capture whole-system provenance in the kernel~\cite{pasquier2017camflow}, partition long-running processes into units~\cite{lee2013beep,ma2017mpi}, or fuse application and audit logs~\cite{hassan2020omegalog,yu2021alchemist}. We build on these ideas; our contribution is their combination for agent tool calls, including subprocesses and non-HTTP Transmission Control Protocol (TCP) connections, with explicit ambiguity and an independent reference.

We compare these efforts along eight dimensions that network attribution requires: \textbf{D1} kernel-level data capture (\eg with eBPF); \textbf{D2} process lineage; \textbf{D3} cgroup membership; \textbf{D4} binding of application context to kernel facts; \textbf{D5} attribution of network activity at connection or flow granularity; \textbf{D6} explicit ambiguity in the output (\ie reporting several candidates rather than forcing a choice); \textbf{D7} validation of attribution against an independent reference; and \textbf{D8} a per-execution management action (\eg accounting or policy). D2 and D3 separate executions with distinct trees or groups. D4 is the only source of distinction inside a shared worker. D5 links traffic, rather than only processes, to executions. D6 and D7 determine whether the output can be trusted for management decisions.

\begin{table}[t]
\centering
\caption{Coverage of attribution dimensions}
\label{tab:dims}
\footnotesize
\setlength{\tabcolsep}{2.6pt}
\begin{tabular*}{\columnwidth}{@{\extracolsep{\fill}}lcccccccc@{}}
\toprule
\textbf{Work} & \textbf{D1} & \textbf{D2} & \textbf{D3} & \textbf{D4} & \textbf{D5} & \textbf{D6} & \textbf{D7} & \textbf{D8} \\
\midrule
AgentSight~\cite{zheng2025agentsight} & \checkmark & \checkmark & $\circ$ & $\circ$ & $\times$ & $\times$ & $\times$ & $\times$ \\
AgentCgroup~\cite{zheng2026agentcgroup} & $\circ$ & $\times$ & \checkmark & $\circ$ & $\times$ & $\times$ & $\times$ & $\circ$ \\
AgentProvenance$^\dagger$~\cite{byteyellow2026agentprovenance} & \checkmark & $\circ$ & \checkmark & \checkmark & $\circ$ & $\circ$ & $\circ$ & $\circ$ \\
ActPlane~\cite{zheng2026actplane} & \checkmark & \checkmark & $\times$ & $\times$ & $\times$ & $\times$ & $\times$ & \checkmark \\
AgentGuard~\cite{shukla2026agentguard} & \checkmark & \checkmark & $\times$ & $\circ$ & $\circ$ & $\times$ & $\times$ & \checkmark \\
ContractWarden~\cite{cui2026contractwarden} & \checkmark & $\circ$ & \checkmark & $\circ$ & $\circ$ & $\times$ & $\times$ & \checkmark \\
OTel eBPF (OBI)~\cite{otel2026obi} & \checkmark & $\circ$ & $\times$ & $\circ$ & $\circ$ & $\times$ & $\times$ & $\circ$ \\
\textbf{\solution{} (this work)} & \checkmark & \checkmark & \checkmark & \checkmark & \checkmark & \checkmark & \checkmark & \checkmark \\
\bottomrule
\end{tabular*}
\scriptsize
\checkmark: present; $\circ$: partial or by other means; $\times$: no evidence in inspected sources; $^\dagger$: also confirmed by running the system (E5ap)
\end{table}

We assessed each system using its published design, documentation, and source code, and marked a dimension only when the inspected sources show the capability (\cf Table~\ref{tab:dims}). Most systems capture kernel events and use process lineage or cgroups (D1--D3), but none reports a candidate set when these signals cannot separate executions; AgentProvenance only marks unmatched events as unresolved (D6). Only AgentProvenance partially validates attribution, through acceptance tests rather than an independent reference (D7). Running it (E5ap, Section~\ref{sec:results}) confirmed its marks: overlapping bindings were resolved to the latest one without reporting ambiguity (D6), and connections of child processes were matched only in part (D2). Network activity is attributed per connection only partially (D5), and the systems that act per execution (D8) key their decisions to process trees, tracked PIDs, or launched tasks rather than to individual tool calls. \solution{} covers all eight dimensions and is the only one that reports ambiguity explicitly and validates attribution against an independent reference.
